%% chapters/00-abstrak-id.tex
%% Abstrak Bahasa Indonesia — 500-800 kata, max 7 kata kunci
%% TIDAK boleh: sitasi, gambar, tabel, persamaan, footnote
%% Aturan: lihat .cursor/rules/12-abstrak-front-matter.mdc

\begin{abstrak}

Prediksi sisa umur pakai (\emph{remaining useful life}, RUL) bantalan
gelinding berperan sentral pada strategi pemeliharaan prediktif modern.
Metode mutakhir berbasis \emph{Long Short-Term Memory} (LSTM) dan
\emph{Transformer} masih menghadapi \emph{trade-off} antara akurasi pada
sekuens panjang, biaya komputasi, dan keterjelasan keputusan model
terhadap fenomena fisis degradasi bantalan. Disertasi ini mengusulkan
\emph{Mamba-xLSTM-Net}, arsitektur hibrida yang mengintegrasikan
\emph{state-space model} selektif (Mamba) untuk pemodelan dinamika
jangka panjang dengan modul \emph{xLSTM} bermemori matriks untuk
mempertahankan presisi rekuren. Arsitektur tersebut dilengkapi modul
interpretabilitas \emph{post-hoc} berbasis \emph{Top-k Sparse
Autoencoder} (SAE) yang dapat dipetakan secara empiris ke frekuensi
karakteristik bantalan (\emph{Ball Pass Frequency Outer race},
\emph{Ball Pass Frequency Inner race}, \emph{Ball Spin Frequency},
\emph{Fundamental Train Frequency}).

Eksperimen dilakukan pada dua dataset publik yaitu PHM2012 dan XJTU-SY
dengan baseline mencakup LSTM, Transformer, Informer, Mamba murni, dan
xLSTM murni. Setiap konfigurasi dilatih dengan lima inisialisasi acak
yang berbeda dan dievaluasi menggunakan metrik \emph{Root Mean Square
Error} (RMSE), \emph{Mean Absolute Error} (MAE), serta \emph{scoring
function} PHM2012. Hasil utama menunjukkan penurunan RMSE rerata
sebesar XX~\% dan MAE sebesar XX~\% dibanding baseline terkuat, dengan
keunggulan paling konsisten pada kondisi operasi beban tinggi. Studi
ablasi mengonfirmasi bahwa kombinasi Mamba dan xLSTM memberikan
kontribusi aditif: penghapusan salah satu modul menyebabkan degradasi
kinerja yang signifikan secara statistik. Analisis interpretabilitas
mengungkap bahwa fitur-fitur laten SAE yang paling sering aktif pada
fase \emph{late-life} berkorespondensi kuat dengan komponen spektral di
sekitar BPFO, BPFI, dan BSF teoretis, memberikan bukti empiris bahwa
model belajar representasi yang konsisten dengan teori klasik diagnosis
getaran bantalan.

Kontribusi utama disertasi ini mencakup: (1) arsitektur hibrida
Mamba-xLSTM-Net yang menggabungkan keunggulan \emph{state-space model}
selektif dan \emph{matrix-memory recurrent}; (2) metodologi
interpretabilitas SAE yang menghubungkan fitur laten dengan konsep
fisis frekuensi karakteristik bantalan; serta (3) bukti empiris
peningkatan kinerja prediksi RUL dengan margin signifikan pada dua
dataset publik yang digunakan secara luas oleh komunitas prognostik.
Hasil disertasi ini membuka jalan bagi pengembangan sistem prognostik
berbasis \emph{deep learning} yang lebih akurat sekaligus dapat
dipercaya oleh para insinyur pemeliharaan di industri.

\vspace{0.5cm}
\noindent
\textbf{Kata kunci:} prediksi sisa umur pakai, bantalan gelinding,
Mamba, xLSTM, sparse autoencoder, interpretabilitas, pemeliharaan
prediktif

\end{abstrak}
